\section{Messung der analogen Signale}

Es wurde zunächst die Signale von Szintillationszähler und Driftkammer direkt an ihren jeweiligen Ausgängen mit einem Oszilloskop betrachtet. Dadurch können die Signalformen der beiden Detektortypen untersucht werden. Die Versorgungshochspannungen werden wie folgt eingestellt: $U_D=\SI{-2.8}{\kilo\V}, U_\text{Szint}=\SI{-1.70}{\kilo\V}$ \cite[2]{515versuchsanleitung}. Während der Durchführung wurde festgestellt, dass bei dieser Einstellung von $U_D$ gelegentlich bereits die Sicherung der Spannungsversorgung ausfiel, vermutlich aufgrund kurzfristig zu hoher Ströme an den Anoden, wenn ein hochenergetisches kosmisches Teilchen in den Detektor einfiel und daher die automatische Überlastungssicherung der Elektronik anging. Die Sicherung musste manuell wieder eingestellt werden, damit die Spannungsversorgung wieder vorhanden war. Insbesondere für die folgende Langzeitmessung war wichtig, diesen Fall zu vermeiden, daher wurde der Wert von $U_D=\SI{-2.8}{\kilo\V}$ als obere Grenze der möglichen Versorgungsspannungen bei der Parametereinstellung betrachtet.

Das Ausgangssignal des Szintillationsdetektors wurde auf einem Oszilloskop als Spannung gegen Zeit aufgenommen.

Das Ausgangssignal eines Anodendrahts der Driftkammer wurde mit einem Tastkopf direkt am äußersten Kontakt der unteren Messkarte (d.h. äußerster Draht der unteren Lage) abgegriffen \cite[2]{515versuchsanleitung}. Damit wurde es nicht durch die folgende Elektronik (Diskriminator) verarbeitet, bevor es gemessen wurde und ist auch unabhängig vom Signal des Szintillationsdetektors. Auch dieses Signal wurde als Spannung-Zeit-Signal mit dem Oszilloskop aufgenommen. Für drei Teilcheneinfall-Ereignisse, die zu einem sichtbaren Detektorsignal der betrachteten Driftkammer-Zelle geführt haben, ist in Abbildung \todo{ref} das aufgenommene Oszilloskopbild dargestellt. Es ist jeweils klar die zu erwartende Signalform eines schnellen Abfalls und einem dazu langsamen, exponentiell verlaufendem Anstieg zurück zur Nulllage der Spannung. Man beachte: Da Elektronen im Gasvolumen ausgelöst werden und sich im Anodendraht Elektronen bewegen, ist die gemessene Spannung negativ. Das Signal hat für alle drei beobachteten Ereignisse eine ähnliche Signaldauer, die aus der Zeitdifferenz zwischen dem Zeitpunkt des Minimums der Spannung und dem Zeitpunkt, bei dem es auf $1/e$ des Minimums wieder angestiegen ist, bestimmt wurde. Der anfängliche Abfall auf das Minimum kann als quasi instantan betrachtet werden bzw, hat keinen signifikanten Anteil an der Dauer des Signalpulses, wie aus den Oszilloskopaufnahmen deutlich wird. Bei die Nutzung eines Diskriminator im finalen Aufbau wird daher nur ein kleiner Teil der sowieso sehr kurzen anfänglichen Abfallsphase bis zum Minimum nicht aufgenommen, der Informationsverlust über den Signalpuls ist also nur klein, während der Vorteil durch Unterdrückung von elektronischem Rauschen in den Daten groß ist. Als Dauer der Signale ist \todo{Werte einsetzen}...

An die aufgenommenen Messwerte werden jeweils ab dem Minimum bis zum Ende des offenbar näherungsweise exponentiellen Wiederanstiegs in die Nulllage eine Exponentialfunktion der Form $U(t)=C\cdot e^{a(t-b)}+b$ angepasst und als Maß der Anpassungsgüte das Bestimmtheitsmaß $R^2$ (geeignet für Verteilungen, die keine Distribution sind \cite{rsquared}) berechnet. Er ergaben sich so Werte von $R^2 > 0,98$, was für eine sehr gute Übereinstimmung der Daten mit dem Modell der Anpassungsfunktion spricht. Die Pulsform ist also als schneller Abfall auf ein Minimum (die ersten Elektronen erreichen die Anode, die sehr nah an ihr ionisiert wurden) und langsamer Abfall (näherungsweise exponential, aufgrund der Ortsverteilung der ausgelösten Elektronen durch das einfallende Teilchen in der Zelle entlang seiner Flugbahn und längeren Driftzeit der entstandenden Ionen im Gas) beschreibbar und entspricht so den Erwartungen an das Signal eines Driftkammer-Detektors \cite[144-145]{leotechniques}.



\jafps{TEK0001}{Oszilloskopaufnahme eines Signals des äußeren, unteren Anodendrahts der Driftkammer in Spannung-Zeit-Darstellung (in orange). Als Anpassungsparameter an die Anpassungsfunktion (in blau) ergaben sich: $a=\SI{-4.843(65)e5}{\per\s}, b=\SI{0.00217(41)}{\V}, c=\SI{-0.0(3.7)e04}{\V}, t_0=\SI{-0.00(34)}{\s}$, als Dauer (grün hinterlegt) $d=\SI{1.8800(28)e-6}{\s}$.}{0.45}
\jafps{TEK0002}{Oszilloskopaufnahme eines weiteren Anodendraht-Signals (in orange). Anpassungsparameter an die Anpassungsfunktion (in blau): $a=\SI{-4.668(61)e+05}{\per\s}, b=\SI{0.00385(45)}{\V}, c=\SI{-0.0(3.6)e+04}{\V}, t_0=\SI{-0.00(32)}{\s}$, als Dauer (grün hinterlegt) $d=\SI{1.8720(28)e-06}{\s}$.}{0.45}
\jafps{TEK0003}{Oszilloskopaufnahme eines weiteren Anodendraht-Signals (in orange). Anpassungsparameter an die Anpassungsfunktion (in blau): $a=\SI{-5.76(13)e+05}{\per\s}, b=\SI{0.00126(52)}{\V}, c=\SI{-0.0(4.8)e+04}{\V}, t_0=\SI{-0.00(40)}{\s}$, als Dauer (grün hinterlegt) $d=\SI{1.3920(28)e-06}{\s}$.}{0.45}


Aufgrund der relativ großen Empfindlichkeit des Szintillationsdetektors wurde die Aufnahme des Szintillatordetektor-Signals im \textit{acquisition}-Modus aufgenommen. \todo{Was heißt das, was macht das mit dem Signal}. So konnte Rauschen im Signal besser unterdrückt werden und das Signal ist besser erkennbar. Das aufgenommene Signal ist in Abbildung \todo{ref} dargestellt.

\jafps{TEK0004}{}{0.5}

%Trotz der längeren Datenaufnahme ist noch ein Rauschen in den Spannungswerten sichtbar, es wurde also aus mehreren Messungen abgeschätzt, dass erst ab einer Spannung von $\SI{-0.550(17)}{\V}$ tatsächlich sichergegangen werden kann, dass es sich um ein Signal eines einfallenden Teilchens handelt.
Das Signal des primären Szintillationsübergangs (für mehr Details zu den weiteren Anregungen im Szintillatormaterial, die Szintillationsphotonen produzieren, siehe bspw. \cite[505-506]{kolanoskiwermerdetektoren}, die die  kleineren Pulsen im Signal erklären), welches als Triggersignal genutzt wird und daher interessant ist, wurde wie oben für die Driftkammersignale auf seine Dauer untersucht. Es ergab sich so eine Pulsdauer von \todo{insert value}. Diese Dauer ist deutlich kürzer als die Signale des Driftkammerdrahts, selbst bei Beachtung der Anstiegszeit des Signals, was hier aufgrund der Einheitlichkeit der Methode nicht präzise durchgeführt wurde. Aus einer groben Abschätzung lässt sich aber sagen, dass die Betrachtung der Anstiegszeit die Dauer des Signals ungefähr verdoppeln würde, was noch immer deutlich kürzer ist, als die Dauer der Driftkammersignale. Dies zeigt, dass der Szintillationsdetektor eine gute Wahl als Trigger für die Aufnahmeelektronik der Driftkammer ist, da er sehr schnell ein Signal anzeigt, sodass die Aufnahme der Driftkammersignale nur eine kurze Verzögerung zum Zeitpunkt des Teilchendurchgangs durch die Kammer hat. Die Signalform entspricht also gut den Erwartungen an den Szintillationsdetektor \cite[505-507]{kolanoskiwermerdetektoren}\cite[21-22,45]{schwandiplom}.
